\subsection{コーシーの平均値の定理}

ロピタルの定理を証明しよう.微分係数の定義に帰着させる方法では,\,$x\to c$の極限で$c$での微分係数だけを使った.今度は$c$と$x$の間のどこかでの導関数の比が必要になる.通常の平均値の定理では,1つの関数の増分を区間の幅で割った.ここでは,2つの関数の増分の比を考える.

$f,g$が$[u,v]$で連続,内部で微分可能で,\,$g'\ne0$なら,ある$\xi\in(u,v)$が存在して
\[
 \frac{f(v)-f(u)}{g(v)-g(u)}=\frac{f'(\xi)}{g'(\xi)}
\]
となる.これをコーシーの平均値の定理という.

この式もロルの定理から導ける.まず,\,$g(v)=g(u)$であればロルの定理からどこかで$g'=0$となってしまうので,分母$g(v)-g(u)$は$0$でない.そこで関数
\[
 H(t)=\{f(v)-f(u)\}\{g(t)-g(u)\}
      -\{g(v)-g(u)\}\{f(t)-f(u)\}
\]
を作ると,\,$H(u)=H(v)=0$である.よって$H'(\xi)=0$となる点$\xi\in(u,v)$があり,
\[
 \{f(v)-f(u)\}g'(\xi)-\{g(v)-g(u)\}f'(\xi)=0
\]
となる.両辺を$\{g(v)-g(u)\}g'(\xi)$で割れば,求める等式を得る.

まず有限の端点$c$で$0$と$0$の形になる場合を考える.端点に$f(c)=g(c)=0$を補えば連続になるので,\,$c$と$x$の間でコーシーの平均値の定理を使える.すると
\[
 \frac{f(x)}{g(x)}=\frac{f'(\xi_x)}{g'(\xi_x)}
\]
となり,\,$\xi_x$は$c$と$x$の間にある.\,$x\to c$なら$\xi_x\to c$だから,右辺の極限が$L$であることから結論が従う.

無限大の形では,もう一段の工夫が要る.例えば$x\to+\infty$で$|g(x)|\to\infty$,\,$L$が実数の場合,\,$\frac{f'(t)}{g'(t)}$が$L$に十分近い範囲の始点$B$を固定する.\,$x>B$に対して
\[
 D_x=\frac{f(x)-f(B)}{g(x)-g(B)}
\]
は,コーシーの平均値の定理により,その範囲のどこかでの$\frac{f'}{g'}$に等しいので,\,$L$に十分近い.元の商は
\[
 \frac{f(x)}{g(x)}
 =D_x\left(1-\frac{g(B)}{g(x)}\right)+\frac{f(B)}{g(x)}
\]
と書ける.\,$B$を固定した後で$x\to+\infty$とすれば,\,$\frac{g(B)}{g(x)}$と$\frac{f(B)}{g(x)}$はそれぞれ$0$へ近づく.そのため元の商も$L$へ近づく.有限の端点で$|g|\to\infty$となる場合も,\,$B$を端点の十分近くに固定して同じ議論を使える.

$0$と$0$の形で端点が$+\infty$の場合は,十分大きい$x<y$の間でコーシーの平均値の定理を使う.増分の比が$L$に近いことを確かめた後,\,$x$を固定して$y\to+\infty$とすると,その比は$\frac{f(x)}{g(x)}$へ近づく.\,$L=+\infty$の場合は「$L$との差が小さい」の代わりに「任意の大きな定数より大きい」を用い,\,$L=-\infty$では不等号を逆にする.同じ議論で,極限が$\pm\infty$の場合も得られる.

\begin{bookpoint}
\textbf{POINT}\par
$0$と$0$,または無限大と無限大の形であることと,分母の導関数が$0$でないことを確かめてから,分子と分母を別々に微分した商の極限を求める.
\end{bookpoint}

\begin{example}{消える次数を比較する}
$\lim_{x\to0}\cfrac{\log(1+x)-x+\cfrac{x^2}{2}}{x^3}$を求めよ.
\end{example}

\paragraphbf{方針}
$0$と$0$の形なので,分子と分母を微分した商を調べる.微分した後に式を整理してから極限をとる.

\paragraphbf{解答}
$\frac13$

\paragraphbf{解法}
分子も分母も$0$に近づき,\,$x\ne0$で分母の導関数$3x^2$は$0$でない.分子を微分すると
\[
 \frac1{1+x}-1+x=\frac{x^2}{1+x}
\]
とまとめられるので
\[
 \cfrac{\left(\log(1+x)-x+\cfrac{x^2}{2}\right)'}{(x^3)'}
 =\frac1{3(1+x)}\longrightarrow\frac13.
\]
この極限が左右どちらからでも$\frac13$なので,ロピタルの定理により求める極限は$\frac13$である.微分後に式を整理すると$x^2$が約分できるため,3回の微分は必要ない.

\paragraphbf{有限テイラー公式による別解}
同じ極限は有限テイラー公式でも求められる.\,\sectionref{節:有限テイラー}の例題と同じく三次の項まで書けば,分子は$\frac{x^3}{3}+R_3(x)$となる.\,$|x|\leq\frac12$では$|f^{(4)}(t)|=\frac6{(1+t)^4}\leq96$なので,\,$|R_3(x)|\leq4|x|^4$と評価でき,\,$x^3$で割った剰余項は$0$へ近づく.どちらの方法も,分子の最初に残る項が三次であることを調べているのだ.

\paragraphbf{条件を省けない例}
ロピタルの定理の条件は省略できない.例えば$x\to0$で$\frac{x^2+1}{x^2}$を考えると,微分した商は$1$だが,元の商は$+\infty$へ発散する.元の分子は$0$にも無限大にも近づかず,定理の仮定を満たしていない.また$x\to+\infty$で
\[
 \frac{x+\sin x}{x}\longrightarrow1
\]
であっても,微分した商$1+\cos x$には極限がない.微分した商の極限が存在しないことから,元の商の極限が存在しないとは結論できない.
